\subsection{Case~5: Collaborative analysis of heart disease cohorts}
\noindent The Cleveland, Hungarian and VA partitions of the UCI Heart Disease dataset were analysed under a common logistic-regression protocol to examine how cross-review informs interpretation across cohorts. One operator managed three roles through separate SciForge accounts and working directories on one Windows device: Role~A analysed Cleveland and coordinated the report, Role~B analysed Hungarian and reviewed the results, and Role~C analysed VA. Methods and aggregate results were shared within the project. Role~B identified a discrepancy between VA test accuracy (75.56\%) and AUROC (0.5629), prompting a revision that considered class composition alongside discrimination. Evaluating the VA test predictions in a follow-up analysis found 35 positives (77.78\%) among 45 VA test cases. The model predicted 38 positives and classified 34 cases correctly, giving 75.56\% accuracy compared with 77.78\% from predicting every case as positive. This follow-up showed that accuracy alone overstated model performance on the imbalanced VA test split. SciForge connected the local analyses, cross-review and report revision within a shared research context (Fig.~\ref{fig:case-heart-disease}; Appendix~\ref{app:case-heart-disease}).
